\documentclass[9pt, twocolumn, a4paper]{ltjsarticle}
\pagestyle{empty}

\usepackage{amsmath, amssymb, amsthm, mathtools}
\usepackage{bm}
\usepackage{physics} 
\usepackage{graphicx}
\usepackage{tikz}
\usepackage{qrcode}
\usepackage[unicode, hidelinks]{hyperref}
\usepackage[margin=15truemm]{geometry}

\title{\vspace{-15mm}\Large\textbf{一端を固定された紐の落下挙動の解析：発表原稿}}
\author{黒瀬 諭一朗（埼玉県立大宮高等学校 物理部）\quad \today}
\date{}

\begin{document}

\maketitle
\section{スライド 1：タイトル}
みなさん，こんにちは．埼玉県立大宮高等学校物理部の黒瀬諭一朗です．本日は「一端を固定された紐の落下挙動の解析」について発表させていただきます．よろしくお願いいたします．

\section{スライド 2：研究の背景と目的}
まず，本研究の背景と目的について説明します．
先行研究における32重振り子の落下シミュレーションにおいて，振り子が初期状態の直線を保ったまま落下しようとし，時間経過とともに固定端側から折れ曲がっていくという興味深い現象が観察されました．
この観察から，「紐という連続的な物体はどのように運動するのか？」という疑問が生じました．そこで本研究では，一端を固定された紐の自由落下挙動を数学的に解明することを目的としました．

\section{スライド 3：解析手法（定式化）}
次に，解析の手法について説明します．
本研究では，紐は「太さ」と「曲げ剛性」を持たない理想的なものと仮定しました．
長さ $l$，質量 $m$ の紐を $N$ 等分し ，長さ $\Delta l = l/N$ の剛体棒と，質量 $\Delta m = m/N$ の質点からなる「$N$ 重振り子」として定式化を行いました．初期状態においては，振り子が鉛直線と角度 $\theta$ をなす直線上で静止しているものとします．

\section{スライド 4：解析手法（アプローチ）}
解析のアプローチです．
落下挙動を調べるにあたり，落下開始直後の時刻 $t$ は微少量であると仮定します．時刻 $t=0$ における各質点の角加速度を $\alpha_k$ と置き，各質点の角度 $\varphi_k$ をテイラー展開を用いて，スライドの式のように近似しました．
解析の流れとしては，運動エネルギーとポテンシャルからラグランジアン $L$ を構成し，初角加速度に関する連立方程式を導出した後 ，$N \to \infty$ の極限を取ることで連続体としての紐の挙動を解析します．

\section{スライド 5〜8：ラグランジアンの構成}
ここからはラグランジアンの構成過程です．
固定端を原点とし，各質点の座標と速度を定義しました．これをもとに系の運動エネルギー $T$ を計算し，同様にポテンシャル・エネルギー $U$ を導出しました．最終的に $L = T - U$ として，スライド8枚目にあるようなラグランジアンを構成しました．

\section{スライド 9：運動方程式の導出}
構成したラグランジアンを用いて，運動方程式を導出します．
ラグランジュ方程式に $t=0$ を代入して計算を行うことで，初角加速度 $\alpha_j$ に関する連立方程式が得られました．

\section{スライド 10：$N \to \infty$ の極限と積分方程式}
次に，連続体としての挙動を解析するため，$N \to \infty$ の極限を取ります．
固定端からの長さを $s$ に対応させて極限を取ると，和の記号が積分に置き換わり，初角加速度 $\alpha(s')$ に関する積分方程式が得られます．この方程式を整理し，両辺を $s$ について微分していくことで，一番下の非常にシンプルな積分方程式を導くことができました．

\section{スライド 11：解の導出とデルタ関数の出現}
得られた積分方程式から解を導出します．
この方程式を $s > 0$ の範囲でさらに微分すると $\alpha(s) = 0$ となります．この条件を満たしつつ元の方程式を成立させるためには，数学的にディラックのデルタ関数 $\delta(s)$ が必要になることが判明しました．
その結果，初角加速度と角度の式には，スライドのようにデルタ関数が出現します．

\section{スライド 12：考察}
この数式的な結果から，どのような物理的挙動が示されているかを考察します．
まず，$s > 0$（固定端以外の自由端側）においては，角度は $\theta$ のままとなります．これは，自由端側の紐が剛体棒のように，初期状態の直線を保ったまま平行移動的に落下しようとすることを示しています．
一方で，固定端である $s = 0$ においては，角度が無限大になるという特異性が現れました．これは「紐全体を直線に保とうとする慣性」と「固定端の拘束力」が衝突した結果であると考えられます．

\section{スライド 13：結論と今後の展望}
最後に，結論と今後の展望です．
本研究により，紐は自由落下を開始した直後，固定端以外の部分は直線を保ったまま落下しようとすることが数学的に確認されました．しかし，固定端にはデルタ関数的な特異性が現れ，「現実には角度が無限大になることはない」という物理的現実との間に矛盾が生じました．
今後の展望としては，現実の紐が持つ「太さ」や「曲げ剛性」を考慮に入れること，そして微小時間だけでなく，有限な時間経過における「折れ曲がりの伝搬」などの落下挙動を解析する手法を構築していきたいと考えています．

以上で発表を終わります．ご清聴ありがとうございました．


\end{document}
